\chapter{General Relativity}
\label{chap:general-relativity}

% Closed question: What reference structure lets distinct observers compare their
% counts as one physical record?

\begin{chapteressay}

A phone reports its position by listening to satellites. Each satellite carries
an atomic clock. Each clock broadcasts a time. The receiver compares arrival
times and solves for position.

The miracle is not merely that the clocks are accurate. The miracle is that the
clocks are corrected. GPS satellites move fast relative to a receiver on
Earth, so special relativity makes their clocks run slow. They are also higher
in Earth's gravitational field, so general relativity makes their clocks run
fast. The net correction is about $38.7\,\mu\mathrm{s}$ per day. Without it,
positions would drift by kilometers per day.

This is the chapter's first lesson: a count is not globally comparable just
because it is precise. It must be carried through a reference structure. A
second on a satellite and a second on Earth are not automatically one physical
record. They become one record when the calibration carries the frame
correction.

The speedometer has the same problem at smaller scale. A wheel sensor counts
rotations. But if the car climbs a mountain road, tires heat, pressure changes,
the road curves, and the wheel's local count no longer tells the whole story by
itself. A GPS speed readout uses a larger frame: satellite clocks, signal
propagation, orbital ephemerides, and relativistic corrections. The local
instrument becomes trustworthy by joining a calibrated reference network.

A second example is the international clock comparison. A national metrology
lab maintains a cesium clock. Another lab maintains another. Each local clock
is excellent, but international time is not simply one local clock promoted to
reality. The clocks are compared, weighted, corrected, and published as a
shared reference. Local counts become a world record through calibration.

The code names this layer with \lean{UniverseTensor}, \lean{REAL}, and
\lean{LOCAL}. \lean{LOCAL} is not a failure; it is a situated claim.
\lean{REAL} is not an oracle; it is what can survive comparison under a
reference structure. \lean{EKG}, living outside the instrument, is the
heartbeat-like reference that lets the comparison happen without exposing its
raw counter as ontology.

The closed question is:

\begin{center}
\emph{What reference structure lets distinct observers compare their counts as
one physical record?}
\end{center}

The answer is the metric and its calibration. General relativity teaches that
time and distance readings are frame-dependent, but not arbitrary. Proper time,
light speed, gravitational redshift, and invariant intervals let different
observers translate their counts into one physical ledger.

Chapter 8 therefore adds reference structure to measurement. Quantum mechanics
showed that records need witnesses. Relativity shows that witnesses need
calibration across frames.

\end{chapteressay}

% Phenomenon arcs use: 1/3 setup -> phenomenon box -> 2/3 structural interpretation.

\section{Calibration Across Frames}
\label{se:calibration-across-frames}

GPS is the modern everyday relativistic instrument. Each satellite carries an
atomic clock and broadcasts its time-of-emission along with its orbital
ephemerides. A receiver on the ground listens to at least four satellites,
notes the time-of-arrival of each signal, computes the propagation delays
implied by the differences, and solves the resulting system for its own
position and the offset of its own internal clock against the satellite
constellation.

The constellation consists of about 24 active satellites in circular orbits
at an altitude of roughly $20{,}200\,\mathrm{km}$ above the Earth's surface,
each travelling at an orbital speed near $3{.}87\,\mathrm{km/s}$. Each
satellite has been calibrated on the ground against the master atomic
standard before launch. After deployment, each satellite continues to be
compared against ground-based standards through downlink telemetry, and its
clock rate is adjusted by command to remain coherent with the constellation.

If the satellite clocks were treated as ordinary Earth clocks, the system
would not work. The discrepancy is large in operational terms even though it
is small in fractional terms, and the calibration was a known requirement
before the constellation was launched.

Two relativistic effects act in opposite directions on each satellite clock.
The first is the slowing predicted by special relativity. A clock moving at
speed $v$ relative to an inertial observer accumulates proper time at the
slower rate $\sqrt{1 - v^2/c^2}$. For a satellite at $v = 3{.}87\,\mathrm{km/s}$,
this works out to a fractional slowing of about $8{.}35 \times 10^{-11}$,
which over one day amounts to $7{.}2\,\mu\mathrm{s}$ lost relative to a clock
at rest on the Earth's surface.

The second is the speeding predicted by general relativity. A clock at higher
gravitational potential (further from a massive body) accumulates proper time
at a faster rate than one deeper in the potential well. The fractional
correction near the Earth's surface is approximately $\Delta\Phi/c^2$, where
$\Delta\Phi = GM_E(1/r_E - 1/r_{\text{orbit}})$ is the gravitational potential
difference between the surface and the orbital altitude. For GPS satellites,
this gives a fractional speedup of about $5{.}28 \times 10^{-10}$, or
$45{.}9\,\mu\mathrm{s}$ per day faster than a surface clock.

The two effects do not cancel. The general relativistic speedup is roughly
six times larger than the special relativistic slowdown, so the net effect
is a satellite clock running fast by approximately $38{.}7\,\mu\mathrm{s}$
per day relative to a clock at the geoid (the Earth's reference surface).
Over the course of a single day, this drift would correspond to a position
error of about $11\,\mathrm{km}$ (since the radio signal travels at
$c \approx 3 \times 10^8\,\mathrm{m/s}$, and $38{.}7 \times 10^{-6} \times
3 \times 10^8 \approx 11{,}600\,\mathrm{m}$). Over a week, the navigation
solution would be useless. Over a year, the satellites would be reporting
positions that bear no relation to any place on Earth.

The fix is built into the satellite clock hardware. Before launch, each
satellite's cesium oscillator is tuned to run slow by exactly the net
relativistic offset. The frequency is offset from the SI value of
$10{.}23\,\mathrm{MHz}$ by a fractional amount of $-4{.}465 \times 10^{-10}$,
so that the apparent frequency seen by a ground-based observer matches the
SI standard once both relativistic effects are accounted for. This factory
calibration is what makes the constellation usable.

The system then carries the calibration certificate through every downstream
operation. The receiver assumes that satellite clocks, after their factory
offset, run at the SI rate as observed from the ground. The propagation
delays, the ionospheric corrections, the satellite ephemerides, and the
receiver's own clock-offset solution all depend on this assumption holding.
If a satellite's frequency drifts from the calibrated rate, the system
detects the drift through redundancy: each receiver solves for its own
position using four satellites; any inconsistency in the resulting solution
indicates that one of the satellites is reporting a time inconsistent with
the others. The ground control segment then issues a correction command, or
in extreme cases removes the satellite from the active constellation.

\begin{phenom}{Einstein Frame Effect}
\label{ph:einstein-frame-effect}

\PhStatement A clock count becomes globally useful only when the reference
frame carrying that clock is included in the calibration.

\PhOrigin Relativity replaced universal coordinate time with frame-dependent
time related by invariant physical structure.

\PhObservation GPS must correct satellite clocks by the combined special and
general relativistic offset before satellite and receiver times can be compared
as one navigation record.

\PhConstraint The measurement may not compare raw clock counts across frames as
though motion and gravitational potential were irrelevant.

\PhConsequence Frame correction is not an optional refinement. It is part of
the calibration certificate that makes the count physical.

\end{phenom}

This is the first place where the word ``gauge'' becomes unavoidable in the
physical book. A local clock reading is a local coordinate. The measurement
is the structure that lets local coordinates be compared. The metric tells
how proper time accumulates along a path. Calibration tells how a particular
instrument enters that structure. Treating two frames' counts as directly
comparable would not be a small error; it would be a category error about
what a count is.

The framework generalizes. Any measurement involving two observers with
different states of motion or different gravitational potentials requires
a frame correction. For most everyday work the correction is too small to
detect --- a clock on the ground floor runs slower than a clock on the tenth
floor by about $10^{-15}$ per meter, undetectable with ordinary instruments.
Modern optical clocks, however, can detect a one-centimeter height difference
through their frequency ratio, and at that precision the floor of a laboratory
becomes a relativistic measurement. Special relativity contributes the
matching half: frames moving relative to each other do not share a single
time coordinate, and the Lorentz transformation maps between them. For
ordinary motions the correction is small but in principle present.

The Lean code formalizes this calibration through the \lean{UniverseTensor}
type. A \lean{UniverseTensor} carries, at each point in spacetime, a metric
tensor specifying how proper time accumulates along admissible paths.
\lean{LOCAL} marks a reading as situated in a specific frame; \lean{REAL}
says when two \lean{LOCAL} readings can be compared after frame correction;
\lean{EKG} holds the heartbeat-like reference scale that lets the comparison
occur.

A reader who has used GPS without thinking about it has been benefiting from
this calibration every day. The phone shows a position correct to a few
meters; those few meters depend on a relativistic correction applied
automatically and continuously by hardware calibrated decades ago. The
structure is invisible because the calibration holds. It would become visible
if it failed.

\section{The Pound--Rebka Experiment}
\label{se:the-pound-rebka-experiment}

In 1959, Robert Pound and Glen Rebka constructed an experiment at the Jefferson
Physical Laboratory at Harvard University to test the prediction of general
relativity that frequency depends on gravitational potential. The experiment
was small: it occupied a single stairwell. The apparatus sent gamma rays from
a source at one end of the tower to an absorber at the other end, separated
by $22{.}6\,\mathrm{m}$ of vertical height. The fractional frequency shift
expected at this height is tiny --- about $2{.}46 \times 10^{-15}$. The
experiment measured this shift to within 10 percent of the predicted value
on its first publication and to within 1 percent of the value in a refined
1964 follow-up.

The smallness of the effect demands a brief calculation. In a uniform
gravitational field with acceleration $g$, the gravitational potential
difference between two points separated by height $h$ is $\Delta\Phi = gh$.
General relativity predicts that a frequency $f$ emitted at the lower
potential and received at the higher one is shifted by
\[
\frac{\Delta f}{f} = -\frac{gh}{c^2}.
\label{eq:chapter8-pound-rebka}
\]
The negative sign indicates a redshift: the received frequency is lower than
the emitted one. For $g = 9{.}81\,\mathrm{m/s^2}$, $h = 22{.}6\,\mathrm{m}$,
and $c = 3 \times 10^8\,\mathrm{m/s}$, the magnitude is $|\Delta f/f| = 2{.}46
\times 10^{-15}$. This is one part in four hundred trillion.

No ordinary frequency standard could detect such a shift directly. The
sharpest atomic resonances available in 1959 had linewidths of about $10^{-7}$
relative to their center frequencies, which is eight orders of magnitude too
broad. To measure a shift of $10^{-15}$, the experimenters needed a frequency
standard whose intrinsic line width was comparable to or narrower than the
shift itself.

The solution was the M\"ossbauer effect, discovered by Rudolf M\"ossbauer in 1958
(and recognized with the Nobel Prize in 1961). The effect is the recoilless
emission and absorption of gamma rays by nuclei bound in a crystal lattice.
When a free nucleus emits a gamma ray, momentum conservation forces the
nucleus to recoil; the recoil takes a small but significant fraction of the
emission energy, broadening the line and shifting its center. When a nucleus
is bound in a rigid crystal lattice, however, the recoil is shared with the
entire lattice. The crystal's enormous effective mass makes the recoil energy
per emission negligible, and the emitted gamma ray has nearly its full
intrinsic line width: extremely narrow.

For the iron-57 isotope used by Pound and Rebka, the relevant transition has
an energy of $14{.}4\,\mathrm{keV}$ and a natural line width of about $5 \times
10^{-9}\,\mathrm{eV}$. The fractional width is $\Gamma/E \approx 3 \times
10^{-13}$. This is still two orders of magnitude broader than the shift to
be measured, but resonance absorption experiments do not require resolving
the line directly; they measure changes in the line's overlap with a
reference, and small shifts produce measurable changes in absorption rate.

The apparatus was the following. At the bottom of the tower, a source
containing iron-57 in an iron matrix emitted gamma rays upward. At the top,
an absorber of identical composition would absorb the gamma rays resonantly
if --- and only if --- the gamma rays arrived at the same energy as the
absorber's transition. If the gamma rays were redshifted by the upward
journey, their energy would no longer match the absorber's resonance, and
absorption would decrease. The detector behind the absorber measured the
rate of transmitted gamma rays; a decrease in transmission indicated
resonance absorption; an increase indicated mismatch.

To restore resonance in the presence of the gravitational redshift, Pound
and Rebka modulated the source's velocity. By moving the source slowly
toward or away from the absorber, they introduced a Doppler shift that
could compensate for the gravitational shift. The compensating velocity is
$v = c \cdot \Delta f / f = c \cdot gh/c^2 = gh/c$. For $g h = 9{.}81
\times 22{.}6 = 221{.}7\,\mathrm{m^2/s^2}$, the compensating velocity is
$v = 221{.}7 / (3 \times 10^8) = 7{.}39 \times 10^{-7}\,\mathrm{m/s}$, or
about $0{.}74\,\mu\mathrm{m/s}$. The source was mounted on a transducer
that could oscillate at this submicrometer-per-second scale.

By scanning the source velocity over a range that included this critical
value, Pound and Rebka mapped the absorption rate as a function of velocity.
The absorption maximum occurred at the velocity that exactly compensated for
the gravitational shift. The measured compensating velocity matched the
predicted $gh/c$ to within the experimental uncertainty.

The 1959 publication reported agreement to about 10 percent. The 1964
refinement, by Pound and Snider, reduced the uncertainty to about 1 percent.
Subsequent improvements with hydrogen masers (Vessot et al., 1976) extended
the precision to better than $10^{-4}$.

\begin{phenom}{Pound-Rebka Effect}
\label{ph:pound-rebka-effect}

\PhStatement A frequency record transported through a gravitational potential
difference must be corrected by the reference structure of spacetime.

\PhOrigin The Pound-Rebka experiment measured gravitational redshift over a
terrestrial height using resonant gamma-ray absorption.

\PhObservation A gamma ray sent upward changes frequency relative to the
receiver by the amount predicted from gravitational time dilation.

\PhConstraint The apparatus may not treat two identical nuclei at different
heights as sharing an uncorrected global frequency standard.

\PhConsequence Even laboratory-scale vertical separation requires metric
calibration when the measurement is precise enough.

\end{phenom}

The experiment is beautiful because it is small. General relativity is not
only about galaxies and black holes. It is about whether two counts can be
compared without lying. Pound-Rebka says: not always. The reference frame
enters the ledger.

The structural significance is the experiment's domesticity. Twenty-two
meters is the height of a modest building. Iron-57 is an ordinary
laboratory isotope. The transducer is a tabletop piece of equipment. The
experiment ran in a stairwell at a university and confirmed general
relativity within 1 percent. The smallness of the setup, combined with
the smallness of the effect, makes the result particularly cleanly
attributable to the gravitational potential difference. There are no
exotic astrophysical environments, no extreme accelerations, no
near-luminal velocities. There is only the ordinary gravitational
potential of the Earth across an ordinary vertical distance.

The experiment is also methodologically clean. The source and absorber are
chemically identical: iron-57 in iron foil. The only relevant physical
difference between them is the gravitational potential at their heights.
Any frequency shift between source emission and absorber resonance can
therefore be attributed to the height difference alone, without competing
explanations involving chemical environment, magnetic field, or
temperature.

For the chapter's framework, the Pound-Rebka result is a calibration
certificate written in the structure of spacetime. The source's frequency
at the bottom of the tower and the absorber's frequency at the top are
not directly comparable as if they were two readings of the same standard.
They are two local readings of different local standards, related by the
gravitational redshift formula. The formula is the calibration; the
factor $gh/c^2$ is the certificate.

The framework extends: every transported frequency record must include
the gravitational potential difference between source and receiver as a
calibration term. For ground-based experiments at modest heights, the
correction is below the resolution of most instruments and can be
ignored. For optical clocks across a city, or for satellite-to-ground
links, the correction is essential. Pound-Rebka established the principle
at the laboratory scale; modern optical clocks and GPS have extended the
practice to engineering applications.

\section{Michelson--Morley}
\label{se:michelson-morley}

Michelson and Morley looked for the ether by comparing light travel times in
perpendicular arms of an interferometer. If Earth moved through an ether, light
should have taken different times along different orientations. Rotating the
apparatus should have shifted the interference fringes.

The expected fringe shift did not appear.

\begin{phenom}{Michelson-Morley Effect}
\label{ph:michelson-morley-effect}

\PhStatement The comparison of light paths is invariant under orientation in a
way that rules out a simple stationary ether frame.

\PhOrigin The 1887 Michelson-Morley experiment searched for Earth's motion
through a luminiferous ether using an interferometer.

\PhObservation Rotating the interferometer produced no ether-wind fringe shift
of the expected size.

\PhConstraint The measurement may not privilege an undetected absolute rest
frame when light-speed comparisons remain invariant.

\PhConsequence The reference structure must be Lorentzian: different observers
relate their counts by transformations preserving the speed of light.

\end{phenom}

This belongs beside GPS and Pound-Rebka because it teaches the other half of
calibration. Gravity changes clock comparison by potential. Relative motion
changes clock and length comparison by Lorentz structure. The old hope was that
one absolute background would make all counts simple. The experiment removed
that hope and replaced it with a better invariant.

Sagnac and Foucault effects fit as neighboring lessons. Rotating frames carry
detectable structure. A pendulum can reveal Earth's rotation. A ring
interferometer can reveal angular velocity. The point is not that all frames
are interchangeable in a lazy sense. The point is that the reference structure
must say exactly how each frame is being compared.

\section{The Schwarzschild Metric}
\label{se:the-schwarzschild-metric}

In January 1916, Karl Schwarzschild, an astrophysicist serving on the Russian
front during World War I, sent Einstein the first exact solution of the
general relativistic field equations. The solution describes the spacetime
outside a static, spherically symmetric mass distribution. In Schwarzschild
coordinates $(t, r, \theta, \phi)$, the line element is
\[
ds^2 = -\left(1 - \frac{2GM}{c^2 r}\right) c^2 dt^2 + \left(1 - \frac{2GM}{c^2 r}\right)^{-1} dr^2 + r^2(d\theta^2 + \sin^2\theta \, d\phi^2).
\]
The quantity
\[
r_s = \frac{2GM}{c^2}
\label{eq:chapter8-schwarzschild-radius}
\]
is the Schwarzschild radius. It depends only on the central mass $M$, the
gravitational constant $G$, and the speed of light $c$. For the Sun's mass,
$r_s \approx 3\,\mathrm{km}$. For the Earth's mass, $r_s \approx
9\,\mathrm{mm}$. For a supermassive black hole at a galactic center with a
mass of $10^9$ solar masses, $r_s \approx 3 \times 10^{12}\,\mathrm{m}$ or
about $20$ astronomical units.

At $r = r_s$, the Schwarzschild coordinate description becomes singular: the
coefficient of $dt^2$ vanishes and the coefficient of $dr^2$ diverges. Early
work (Hilbert, Droste, others) treated this as a physical singularity.
Beginning with Lema\^itre in 1933 and clarified by Kruskal and Szekeres in
1960, it became understood that the singularity at $r = r_s$ is a coordinate
artifact: it appears in Schwarzschild coordinates but not in better-chosen
coordinates (such as Eddington-Finkelstein or Kruskal-Szekeres coordinates).
The geometry at $r = r_s$ is perfectly regular; what fails is the
Schwarzschild coordinate chart's ability to cover it smoothly.

What does not fail is the causal structure. The surface $r = r_s$ is a
\emph{causal horizon}: a hypersurface across which signals can travel inward
but not outward. Light cones tilt as $r$ decreases; at $r = r_s$, the
outward light cone is exactly tangent to the horizon, meaning that outgoing
light rays neither escape nor fall in but stay at $r_s$. For $r < r_s$, even
light rays initially directed outward must travel inward; nothing in the
interior can send a signal that crosses $r_s$ from inside to outside. The
horizon is a one-way membrane.

This horizon is the chapter's most severe instance of frame-dependent
admissibility. Two observers --- one distant, one falling in --- can both be
measuring the same falling object, and their two descriptions are both
locally admissible, but the two descriptions cannot be combined into a single
globally communicable record.

Consider a clock falling radially into a Schwarzschild black hole. From the
perspective of a distant observer who watches via light signals, the clock's
ticks become increasingly redshifted as it approaches the horizon. The
fractional redshift grows as $(1 - r_s/r)^{-1/2}$, which diverges as $r
\to r_s$. The distant observer sees the clock's tick rate slow to zero and
the clock's apparent position freeze at the horizon. Light from the clock
takes longer and longer to escape, and after some finite emission time, no
further light can escape at all. From the distant observer's point of view,
the clock asymptotically approaches the horizon but never crosses it.

From the perspective of the clock itself, however, nothing dramatic happens
at the horizon. The clock falls smoothly through $r = r_s$ in finite
proper time, as measured by the clock's own ticking. For a clock falling
from rest at infinity into a Schwarzschild black hole of mass $M$, the
proper time from horizon crossing to the central singularity at $r = 0$ is
$\tau = (\pi/2)(r_s/c)$, which for a stellar-mass black hole is on the
order of microseconds. The clock crosses the horizon, ticks a few more
times, and ends at the central singularity. The clock's own record contains
the horizon crossing as a perfectly ordinary event.

The two records are both correct in their own frames, but they cannot be
merged. The distant observer's record contains no event that corresponds to
horizon crossing; the falling clock's record contains the horizon crossing
in finite proper time. There is no third, neutral coordinate system in which
both records can be expressed as readings of one global clock. The
Schwarzschild geometry simply does not admit such a coordinate system.

This is the chapter's structural lesson stated most starkly. Two locally
admissible measurements may not be combinable into one global record if the
spacetime geometry prevents the causal exchange required for the
combination. The geometry is not an obstruction to physics; it is what
physics says about admissibility. The horizon's existence is what
prevents the merging, not a failure of measurement technique.

\begin{phenom}{Event Horizon Effect}
\label{ph:event-horizon-effect}

\PhStatement A reference structure can make two locally admissible descriptions
incompatible as one outward communicable record.

\PhOrigin Schwarzschild geometry revealed a horizon in the spacetime outside a
compact mass, later understood as the causal boundary of a black hole.

\PhObservation To a distant observer, infalling matter appears increasingly
redshifted near the horizon; for the infalling observer, crossing the horizon
can occur in finite proper time.

\PhConstraint The measurement may not demand one global coordinate story where
the causal structure permits only frame-dependent descriptions.

\PhConsequence General relativity separates local admissibility from global
communicability. A record can be real locally and unavailable outward.

\end{phenom}

The event horizon is a severe version of a modest lesson. Counts are local.
Comparison requires transport. If transport cannot carry a signal back, the
outside ledger cannot include the inside event as an ordinary returned report.

The Schwarzschild geometry has been experimentally confirmed in several
contexts beyond black holes. The orbit of Mercury exhibits a precession of
about $43$ arcseconds per century that Newtonian gravity cannot explain;
general relativity, applied to the Schwarzschild geometry around the Sun,
predicts the precession to within the measurement uncertainty. The bending
of starlight passing near the Sun, measured during the 1919 solar eclipse
by Eddington's expedition, matched the Schwarzschild prediction of $1{.}75$
arcseconds at the solar limb to within experimental error. The Shapiro
delay --- the additional travel time of light passing near a massive body
--- has been measured in the solar system to one part in $10^5$ using
spacecraft radar.

The Event Horizon Telescope's images of the supermassive black hole at the
center of M87 (2019) and of Sagittarius A* (2022) show the shadow cast by
the horizon at the predicted angular size. The shadow's diameter, the ring
of light around it, and the asymmetry produced by the black hole's rotation
all match the predictions of the Schwarzschild and Kerr geometries.

The horizon's existence is therefore not a theoretical novelty; it is an
observed feature of nature. The chapter's claim is structural: the
existence of horizons is forced by the geometry, and the geometry is
forced by the requirement that admissible measurements be self-consistent.
The frame-dependent descriptions are not artifacts of choosing the wrong
coordinates; they are the structural form that admissible measurement
takes in curved spacetime.

This is why the chapter talks about reference structure instead of a
God's-eye view. The metric is not a decorative field painted over space.
It tells which counts can be compared, how clocks accumulate proper time,
and where causal communication closes. The Lean code's
\lean{UniverseTensor} formalizes this: it carries the metric structure
that determines admissibility of comparisons; the \lean{LOCAL} certificate
marks readings as situated in specific frames; the \lean{REAL} relation
indicates when two \lean{LOCAL} readings can be merged into a global
record. For Schwarzschild geometry inside the horizon, \lean{REAL} fails
to admit the merge: there is no transport from inside the horizon to a
distant observer's frame. The code's discipline aligns with the geometry's
discipline.

\section{The Atomic Clock as EKG}
\label{se:the-atomic-clock-as-ekg}

The SI second is defined by a cesium-133 transition: $9,192,631,770$ periods of
radiation associated with the hyperfine splitting of the ground state. That
number is not a philosophical foundation. It is a calibration choice made
physical by reproducible atomic behavior.

An atomic clock is therefore an EKG in the book's sense. It is a heartbeat-like
reference process against which other counts can be compared. The heartbeat is
not exposed as the thing being measured. It is held as the reference that lets
measurement proceed.

\begin{phenom}{Atomic EKG Effect}
\label{ph:atomic-ekg-effect}

\PhStatement A shared time record requires a repeatable physical heartbeat that
can act as a reference without becoming the measured object itself.

\PhOrigin Atomic time standards use stable atomic transitions to define and
compare seconds across laboratories and instruments.

\PhObservation Cesium clocks and related standards let local laboratories
compare time by referring counts to a reproducible transition frequency.

\PhConstraint The reference heartbeat must be carried as calibration; exposing
its raw counter as ontology confuses the scale with the measured event.

\PhConsequence Physical timekeeping is a calibrated comparison of heartbeats,
not direct access to an unframed universal clock.

\end{phenom}

This is the bridge from relativity back into the project's calibration work.
\lean{EKG} holds a reference. \lean{REAL} requires comparison across witnessed
truth. \lean{LOCAL} keeps frame-dependent variation honest. General relativity
does the same in physics: it lets local counts become a shared record without
pretending locality disappeared.

\begin{bridge}

The chapter's question was what reference structure lets distinct observers
compare their counts as one physical record.

The answer is calibrated metric structure. GPS requires relativistic clock
corrections. Pound-Rebka measures gravitational frequency shift. Michelson-
Morley removes the ether and leaves Lorentz invariance. Schwarzschild geometry
shows that communication itself can have a boundary. Atomic clocks provide the
heartbeat by which local counts are compared.

The next question is what happens when transport around a loop fails to return
unchanged. Chapter 9 asks what records the failure of transport to commute.

\end{bridge}

\begin{coda}{An International Time-Standard Conference}

The conference room is full of clock people. Each institute has
brought data from its own laboratory. Every clock has been cared
for locally. None of them is enough by itself.

The first rule of the conference is humility. A local clock is not
wrong because it is local; it is incomplete as a world standard.
Its tick is excellent inside its own laboratory. To become part of
international time, it must be compared.

The comparison is not casual. Data arrive with uncertainty
estimates. Some labs receive more weight than others. The
committee does not crown the loudest clock. It constructs a shared
record.

This performs the chapter. The laboratory clocks are \lean{LOCAL}. Their
measurements are not discarded. They are transported through calibration. The
published international timescale is \lean{REAL} in the book's restrained
sense: not metaphysical authority, but the record that survives the comparison
protocol. The atomic transitions are \lean{EKG}, the heartbeat references. The
matrix of comparisons is the human version of \lean{UniverseTensor}.

A clock at high altitude must be corrected differently from a clock at sea
level. A satellite link must include propagation delay. A lab with a better
uncertainty budget may receive more weight. None of this embarrasses time. It
is how time becomes measurable across the world.

At the end, the conference publishes corrections. Each institute can now say
how its local clock relates to the shared scale. The result is not one clock
dominating all others. It is a calibrated agreement among many clocks, each
still local, each now comparable.

General relativity has the same temperament. It does not abolish local
readings. It teaches the discipline by which local readings can meet. The
world record is not raw sameness. It is calibrated comparison.

When everyone leaves the room, the clocks keep ticking in their separate
laboratories. The shared second lives in the work that lets those ticks agree.

\end{coda}
